\chapter{Geometría}

En esta sección exploraremos lo básico para resolver
problemas de geometría de olimpiadas.

Esta sección no tiene como intención crear una
construcción axiomática de los factos que se exponen,
y se asume un conocimiento previo
e intuición mínima (muy mínima).

\section{Ángulos y triángulos}
Lo más básico que se puede discutir acerca
de los ángulos es cómo interactuan con
rectas paralelas, y es algo que probablemente
ya conozcas:
\begin{theorem}
	Una recta forma un mismo ángulo
	con otras dos rectas si éstas son
	paralelas.

	\input{figures/paralelas.tex}
\end{theorem}

A la igualdad entre los ángulos marcados
en rojo se los llama ángulos correspondientes
y a la igualdad entre el ángulo rojo
de abajo y el azul se le denomina como alternos internos.

Y resulta ser que esto se puede usar para
probar otra cosa que probablemente también
conozcas:

\begin{proposition}\label{proposition:sumangle180}
	Los ángulos internos de un triángulo suman $180^\circ$.
\end{proposition}

Pensándolo un poco, resulta ser bastante sorprendente
que, sin importar su forma, cualquier triángulo cumple
que la suma de sus ángulos internos es $180^\circ$.
La manera de probarlo es la siguiente:

\vspace{0.5cm}
\centering
\begin{asy}
	size(6cm);
	pair A = dir(110);
	pair B = dir(210);
	pair C = dir(330);

	draw((A - (0.5, 0))--(A + (1, 0)));
	draw(A--B--C--cycle);
	dot("$A$", A, (-0.5, 1));
	dot("$B$", B, SW);
	dot("$C$", C, SE);
	label("$\ell$", (0.7, 1));

	markangle(A, C, B, radius=10, red);
	markangle(C, B, A, radius=10, blue);
	markangle((A - (1, 0)), A, B, radius=10, blue);
	markangle(C, A, (A + (1, 0)), radius=10, red);
\end{asy}
\par
\raggedright

Seleccionamos un vértice y dibujamos la recta $\ell$
paralela al lado opuesto que pasa por ese vértice,
por alternos internos vemos que $\dang(\ell, AB) = \angle B$
y $\dang(AC, \ell) = \angle C$
y resulta ser que $\angle A + \angle B + \angle C = 180^\circ$
ya que juntos forman una recta (la recta $\ell$).

\begin{remark}
	Aquí la notación $\dang(\ell_1, \ell_2)$
	simboliza el ángulo por el que debe rotarse la recta $\ell_1$
	alrededor del punto de intersección
	en sentido antihorario para que coincida con la
	recta $\ell_2$.
\end{remark}

\subsection*{Problemas para esta sección}
\begin{problem}
Demuestre que la suma de los ángulos internos
de un polígono
de $n$ lados que no se intersecan
es $180(n - 2)^\circ$.
\addhints{hint:sumapol1, hint:sumapol3}
% , hint:sumapol2
\end{problem}

\section{Razones y Semejanza}
Las razones en geometría son de extrema
importancia, consisten en hablar
de las proporciones que tienen ciertos
segmentos o áreas, y convenientemente existen
muchos resultados que nos ayudan a trabajar
con éstos, uno de ellos es
el Teorema de Thales, que lo enunciamos de la siguiente
manera:
\begin{theorem}[Teorema de Thales]
	Sea $ABC$ un triángulo y sean $X$ e $Y$
	puntos en las rectas $AB$ y $AC$,
	entonces las rectas $XY$ y $BC$
	son paralelas si y solo si
	\[\frac{AX}{AB} = \frac{AY}{AC}.\]
\end{theorem}

\centering
\begin{asy}
	size(6cm);
	pair A = dir(110);
	pair B = dir(210);
	pair C = dir(330);
	pair X = midpoint(A--B);
	pair Y = midpoint(A--C);

	draw(A--B--C--cycle);
	draw(X--Y, dashed);
	dot("$A$", A, NW);
	dot("$B$", B, SW);
	dot("$C$", C, SE);
	dot("$X$", X, W);
	dot("$Y$", Y, E);
\end{asy}
\par
\raggedright

Una demostración mejor de lo que creo
soy capaz de explicar se encuentra
\href{https://www.youtube.com/watch?v=d7QXlL_Ztvc}{aquí}.

\begin{example}
	Sea $ABCD$ un cuadrilátero.
	Se traza por $C$ la recta paralela
	a $AD$ e interseca a $BD$ en $M$.
	De manera similar se traza la recta paralela
	a $BC$ por $D$ e interseca a $AC$
	en $N$. Demuestre que $MN$ es paralela a $AB$.
\end{example}

Al principio este problema puede dar miedo,
pero nos damos cuenta de que si $P$ es
la intersección de las diagonales $AC$
y $BD$, entonces, por el teorema de
Thales es suficiente probar que
$\frac{PM}{PB} = \frac{PN}{PA}$.
Por suerte tenemos paralelas
que involucran justamente a esos segmentos.
Sin miedo al éxito aplicamos Thales
con $P$ y $BC \parallel ND$ y obtenemos:
\[\frac{PD}{PB} = \frac{PN}{PC},\]
y de manera similar con $P$ y $AD \parallel MC$:
\[ \frac{PM}{PD}= \frac{PC}{PA},\]
para finanizar multiplicamos ambas
ecuaciones y concluimos que
\[\frac{PM}{PB} = \frac{PN}{PA},\]
que es justo lo que queríamos probar.

\centering
\begin{asy}
	size(6cm);
	pair A = dir(110);
	pair B = dir(210);
	pair C = dir(330);
	pair D = dir(40);

	pair aux = A - D + C;

	pair M = extension(B, D, C, aux);
	aux = B - C + D;
	pair N = extension(A, C, D, aux);

	pair P = extension(A, C, B, D);

	draw(A--B--C--D--cycle);
	draw(A--C);
	draw(B--D);
	draw(C--M);
	draw(D--N);
	draw(M--N, dashed);

	dot("$A$", A, N);
	dot("$B$", B, SW);
	dot("$C$", C, SE);
	dot("$D$", D, NE);
	dot("$M$", M, S);
	dot("$N$", N, NE);
	dot("$P$", P, (0.36, 1));
\end{asy}
\par
\raggedright

Podemos analizar más propiedades relacionadas a
la forma de los triángulos,
sin ir más allá, nos es difícil
imaginar dos triángulos, uno más grande que
otro pero ambos con la misma forma, al
tener la misma forma entonces también tienen
los mismos ángulos internos, las mismas razones,
¿cómo podemos hacer uso de esto en los problemas?,
simple: si dos triángulos tienen
la misma ``forma'' pero
con distinto (o posiblemente igual) tamaño decimos que son
\textit{triángulos semejantes}.

Aquí van los tres criterios de semejanza principales:
\begin{theorem}[Criterio AAA, Ángulo-Ángulo-Ángulo]
	Dos triángulos $ABC$ y $XYZ$ son semejantes
	si y solo si $\angle A = \angle X$, $\angle B = \angle Y$ y
	$\angle C = \angle Z$. Más aún,
	si dos de esas igualdades de ángulos son ciertas,
	entonces también lo es la tercera.
\end{theorem}

\begin{remark}
	La aclaración se desmitifica sabiendo que
	la suma de los ángulos internos de un triángulo es
	constante (igual a $180^\circ$), ya que si
	se conocen dos ángulos, entonces es posible encontrar el tercero.
\end{remark}

\begin{theorem}[Criterio LLL, Lado-Lado-Lado]
	Dos triángulos $ABC$ y $XYZ$ son semejantes si y solo si
	\[\frac{AB}{XY} = \frac{BC}{YZ} = \frac{CA}{ZX}.\]
\end{theorem}

\begin{theorem}[Criterio LAL, Lado-Ángulo-Lado]
	Dos triángulos $ABC$ y $XYZ$ son semejantes si y solo si
	$\angle A = \angle X$ y
	\[\frac{AB}{XY} = \frac{AC}{XZ}.\]
\end{theorem}

Este último criterio resulta ser problemático
en algunas ocasiones porque el orden Lado-Ángulo-Lado
es muy importante, si en su lugar tuvieramos que
$\angle A = \angle X$ y $\frac{AB}{XY} = \frac{BC}{YZ}$,
entonces no podemos afirmar que los triángulos
sean semejantes.

\begin{question*}
	¿Por qué pasa esto?
\end{question*}

\begin{definition}[Triángulos directa y opuestamente semejantes]
	Dos triángulos son directamente semejantes si
	el listado de sus vértices correspondientes
	están ambos en el mismo sentido de
	las agujas del reloj y en caso contrario
	son opuestamente semejantes.
\end{definition}

A esto también se lo puede ver como que
si es posible cambiar el tamaño de uno de
los triángulos y trasladarlo al lugar del otro
para que coincida con este entonces son directamente
semejantes pero si es necesario \textit{reflejar}
uno de los triángulos para obtener esto
entonces son opuestamente semejantes.

No establezco una definición de semejanza
entre dos triángulos porque todas las
anteriores son equivalentes.

\begin{lemma}\label{lemma:rotomotecia}
	Sean $OAB$ y $OXY$ dos triángulos \textit{directamente}
	semejantes, entonces $OAX$ y $OBY$
	son directamente semejantes.
\end{lemma}

Hacemos uso del criterio LAL. La semejanza
nos dice que $\angle AOB = \angle XOY$,
pero como los ángulos $\angle AOX$
y $BOY$ comparten el ángulo
$BOX$, entonces éstos son iguales.
Por la semejanza también tenemos
\[\frac{OA}{OX} = \frac{OB}{OY},\]
reordenando la razón vemos que
\[\frac{OA}{OB} = \frac{OX}{OY},\]
pero esto junto con la igualdad de
ángulos completa el criterio LAL
y por ende $OAX$ y $OBY$ son semejantes.

\centering
\begin{asy}
	size(8cm);
	pair O = (0, 0);
	pair A = (-2, -4);
	pair B = (-1, -6);

	picture p1;
	picture p2;

	pair X = (rotate(40)*A)*0.8;
	pair Y = (rotate(40)*B)*0.8;

	draw(p1, O--A--B--cycle);
	draw(p1, O--X--Y--cycle);

	filldraw(p1, O--A--X--cycle, lightblue + opacity(0.3));
	filldraw(p1, O--B--Y--cycle, lightblue + opacity(0.3));


	dot(p1, "$O$", O, N);
	dot(p1, "$A$", A, W);
	dot(p1, "$B$", B, SW);
	dot(p1, "$X$", X, NE);
	dot(p1, "$Y$", Y, E);

	filldraw(p2, O--A--B--cycle, lightblue + opacity(0.3));
	filldraw(p2, O--X--Y--cycle, lightblue + opacity(0.3));

	draw(p2, O--A--X--cycle);
	draw(p2, O--B--Y--cycle);

	dot(p2, "$O$", O, N);
	dot(p2, "$A$", A, W);
	dot(p2, "$B$", B, SW);
	dot(p2, "$X$", X, NE);
	dot(p2, "$Y$", Y, E);

	add(p1);
	add(shift(7, 0)*p2);
\end{asy}
\par
\raggedright

\begin{example}
	Sea $ABC$ un triángulo y $M$, $N$ y $P$
	los puntos medios de los lados $AB$, $BC$ y $AC$
	respectivamente. Demuestre que $NPM$
	es semejante a $ABC$.
\end{example}
Aquí hacemos uso del teorema de Thales
y de un criterio de semejanza.
Por ser $M$ y $N$ puntos medios de $AB$ y $BC$,
entonces
\[\frac{BM}{BA}=\frac{BN}{BC} = \half\]
pero por el teorema de Thales esto implica
que $MN$ es paralela a $AC$. De manera
análoga sabemos también que $MP\parallel BC$
y $NP\parallel AB$.
Por las paralelas sabemos que $\angle NMP = \angle APM = \angle C$
de manera completamente análoga concluimos que
$\angle MPN = \angle B$
y $\angle PNM = \angle A$, concluyendo
la demostración por el criterio AAA.

\centering
\begin{asy}
	size(6cm);
	pair A = dir(110);
	pair B = dir(210);
	pair C = dir(330);

	pair M = midpoint(A--B);
	pair N = midpoint(C--B);
	pair P = midpoint(A--C);

	draw(A--B--C--cycle);
	draw(M--N--P--cycle);
	dot("$A$", A, NW);
	dot("$M$", M, W);
	dot("$B$", B, SW);
	dot("$N$", N, S);
	dot("$C$", C, SE);
	dot("$P$", P, E);

	markangle(B, A, C, radius=10, red);
	markangle(B, M, N, radius=10, red);
	markangle(P, N, M, radius=10, red);

	markangle(C, B, A, radius=10, blue);
	markangle(C, N, P, radius=10, blue);
	markangle(M, P, N, radius=10, blue);

	markangle(N, M, P, radius=10, green);
	markangle(A, P, M, radius=10, green);
	markangle(A, C, B, radius=10, green);

\end{asy}
\par
\raggedright

\begin{remark}
	A $MNP$ se lo conoce como el triángulo
	medial de $ABC$.
\end{remark}

\subsection*{Problemas para esta sección}
\begin{problem}
Muestre que el criterio Lado-Lado-Ángulo
sí funciona si el ángulo en cuestión
es igual a $90^\circ$.
\end{problem}

\begin{problem}
Muestre que el \autoref{lemma:rotomotecia}
no es cierto en general
si no se restringe que los triángulos
sean directamente semejantes.
\end{problem}


\section{Circunferencias}
Empezamos con qué es una
circunferencia.

\begin{definition}
	Una circunferencia con centro $O$ y radio
	$r$ es el espacio geométrico de
	todos los puntos que distan $r$ de $O$.
\end{definition}

Las circunferencias son fundamentales
por que en particular si $A$
y $B$ son dos puntos en la circunferencia,
entonces $OA = OB$, lo que nos da un triángulo
isósceles de gratis además de $\angle OAB = \angle OBA$.

Este es un hecho básico: %TODO: Reparafrasear esa basura

\begin{theorem}
	Para cualesquiera tres puntos en el
	plano, no todos colineales, existe
	una única circunferencia que pasa
	por esos tres puntos.
\end{theorem}

Una manera de demostrar esto sería
probar que existe un punto $O$ en el plano
tal que $OA = OB = OC$, de esta manera
la circunferencia de centro $O$ y radio $OA$
pasaría por los tres puntos, que es lo que estamos buscando.

Recordamos que esto es cierto. %TODO: Parafrasear esto.
\begin{lemma}
	La mediatriz de un segmento $XY$ es
	la recta perpendicular a
	$XY$ que pasa por su punto medio.
	Esa recta cumple que para todo punto $P$
	en ella se cumple que $PX = PY$.
	Además, si $Q$ es un punto tal que $QX = QY$,
	entonces $Q$ está en la mediatriz.
\end{lemma}

No demostraremos esto por ser muy elemental.
Con este lema en cuenta
vemos que si trazamos la mediatriz
de, por ejemplo, el segmento $AB$,
entonces cualquier punto $P$
en esta cumple que $PA = PB$,
si tan solo pudiésemos encontrar un
punto en esa mediatriz tal que $PA = PC$,
entonces ese punto sería el punto $O$
que estamos buscando. Pero se nos ocurre que
la mediatriz del segmento $AC$
cumple que cualquier punto en esta equidista
de $A$ y $C$, por lo que si definimos a
$O$ como la intersección de las dos mediatrices,
entonces como está en la mediatriz de $AB$,
tenemos $OB = OA$,
pero como está en la mediatriz de $AC$, entonces
$OA = OC$, en conclusión:
\[OA = OB = OC,\]
que es el punto que estábamos buscando.
Más aún, como $OB = OC$, entonces $O$ también
está en la mediatriz de $BC$. De esta manera
deducimos un resultado más fuerte:

\begin{theorem}
	Si $ABC$ es un triángulo, entonces
	las mediatrices de los lados $AB$,
	$BC$ y $AC$ intersecan en el centro
	de la circunferencia que pasa por
	$A$, $B$ y $C$.
\end{theorem}

A este punto $O$ de lo llama el
\textit{circuncentro} del triángulo $ABC$.
Decimos \textbf{la} porque esa circunferencia
es única.

\begin{question*}
	¿Por qué existe una única circunferencia
	que pasa por tres puntos en el plano?
\end{question*}

Ahora presentamos lo que es, quizás,
el resultado más sorprendente relacionado
a las circunferencias:

\begin{theorem}[Medida del ángulo central]
	Sean $A$, $B$ y $C$ tres puntos en una circunferencia
	de centro $O$ con $A$ del mismo lado
	de la recta $BC$ que $O$,
	entonces $\angle BOC = 2\angle BAC$.
\end{theorem}

\begin{figure}[h]
	\begin{center}
		\input{figures/angulo_central.tex}
	\end{center}
	\caption{El ángulo central es el doble del arco inscrito.}
\end{figure}

Para probar esto no es necesario nada
más y nada menos que con los triángulos isósceles
y las igualdades de ángulos que nos dan.
Dibujamos el rayo $AO$,
claramente $OAB$ y $OAC$ son triángulos isósceles,
por ende $\angle OBA = \angle OAB = \alpha$ y $\angle OCA = \angle O = \beta$,
en particular por \autoref{proposition:sumangle180}
tenemos $\angle BOA = 180 - 2\alpha$
y $\angle COA = 180 - 2\beta$.
Y como $\angle BOA + \angle COA + \angle BOC = 360$
concluimos
\[\angle BOC = 360 - (180 - 2\alpha - 2\beta) = 2\alpha + 2\beta = 2\angle BAC.\]

Esto implica que si $ABCD$ es un cuadrilátero convexo
inscrito en una circunferencia de centro $O$,
entonces $\angle BAC = \angle BDC$ por que
ambos serán iguales al ángulo central $\angle BOC$,
a esto se lo denomina arco capaz.

\begin{figure}[H]
	\begin{center}
		\input{figures/arco_capaz.tex}
	\end{center}
	\caption{Arcos capaces.}
\end{figure}

Podemos deducir más resultados respecto
a las circunferencias, como lo es el
siguiente:
\begin{theorem}
	Sea $ABCD$ un cuadrilátero convexo inscrito
	en una circunferencia, entonces
	\[\angle BAD + \angle BCD = 180^\circ.\]
\end{theorem}

Esto no es difícil de deducir por
la medida del ángulo central. Sea
$O$ el centro de la circunferencia,
entonces
\[\angle BAD + \angle BCD = \half\angle BOD + \half\angle DOB = \half\left(\angle BOD + \angle DOB\right) = \half 360^\circ = 180^\circ.\]
\begin{figure}[H]
	\begin{center}
		\begin{asy}
	size(6cm);
	pair A = dir(80);
	pair B = dir(160);
	pair C = dir(210);
	pair D = dir(330);

	pair O = (0, 0);

	markangle(B, A, D, radius=10);
	markangle(D, C, B, radius=10);

	draw(B--O--D);
	draw(A--B--C--D--cycle);
	draw(unitcircle);

	dot("$A$", A, NE);
	dot("$B$", B, NW);
	dot("$C$", C, SW);
	dot("$D$", D, SE);
	dot("$O$", O, NE);
\end{asy}

	\end{center}
	\caption{Ángulos opuestos suman $180^\circ$.}
\end{figure}

\begin{lemma}\label{lemma:righttrianglemindpoint}
	Sea $ABC$ un triángulo recto en $B$
	y sea $O$ un punto en el segmento
	$AC$. Si dos de las logitudes
	$OA$, $OB$ y $OC$ son iguales,
	entonces también lo es la tercera.
\end{lemma}

Una vez, durante un examen de ingreso estábamos
vendiendo panchos y ofrecíamos uno gratis si
eran capaces de resolver este sencillo problema
(nadie lo resolvió):

\begin{example}
	Halle el valor de $x$ que se indica en la figura.
	\begin{figure}[H]
		\begin{center}
			\begin{asy}
	size(6cm);
	pair D = dir(30);
	pair C = dir(-50);
	pair A = dir(90);
	pair B = dir(190);

	pair P = extension(A, C, B, D);

	draw(A--B--C--D--cycle);
	draw(A--C);
	draw(B--D);
	markrightangle(A, P, B);

	dot("$A$", A, NW);
	dot("$B$", B, SW);
	dot("$C$", C, S);
	dot("$D$", D, NE);
	dot(P);

	label("$40^\circ$", A, (2.9, -3.9));
	label("$40^\circ$", B, (3, 0));
	label("$30^\circ$", C, (-0.5, 9));
	markangle(P, B, A);
	label("$x$", B, (8, 6.5));
\end{asy}


		\end{center}
	\end{figure}
\end{example}

A primera vista puede parecer que
con saber que la suma de los ángulos
interiores de un triángulo
suman $180^\circ$ puede ser suficiente
para hallar $x$, pero
resulta ser que no,
y para ello necesitamos una
herramienta que vimos antes:
la semejanza.

Sea $P$ la intersección de las diagonales,
entonces los triángulos $PAD$
y $PBC$ son opuestamente semejantes
por $AAA$, esto implica que
\[\frac{PA}{PB} = \frac{PD}{PC},\]
reordenando la ecuación obtenemos
\[\frac{PA}{PD} = \frac{PB}{PC}.\]
Como $\angle APB = \angle CPD$
por ángulos opuestos, entonces por
el criterio LAL, los triángulos $PAB$
y $PDC$ son semejantes,
lo que implica que
\[x = \angle PBA = \angle DCP = 30^\circ.\]

Resulta que este razonamiento puede
ser aplicado para deducir un resultado más general:
\begin{problem}
Sea $ABCD$ un cuadrilátero convexo
con $\angle CAD = \angle CBD$,
demuestre que $\angle DBA = \angle DCA$.
\end{problem}

% TODO: DESARROLLAR el lema
\begin{lemma}
	Sea $H$ el ortocentro de un triángulo $ABC$, entonces las
	refexiones de $H$ sobre el lado $BC$ y sobre el punto medio de $BC$
	respectivamente se encuentran en el circuncírculo de $ABC$.
	Y $A[reflexion de H sobre midpoint de BC]$ es diámetro de
	$ABC$.
\end{lemma}

\section{Cuadriláteros cíclicos}
Este término es quizás el más común y
útil en geometría de olimpiada.

La palabra ``cíclico'' pueda dar miedo,
aunque es en realidad bastante sencillo.

\begin{definition}[Cuadrilátero cíclico]
	Un cuadrilátero es cíclico si existe
	una circunferencia que pasa por sus cuatro vértices.
\end{definition}

Al principio puede parecer que si sabemos
que un cuadrilátero es cíclico no nos
da mucha información, pero esto no es
así ya que nos da una infinidad de
igualdades de ángulos y semejanzas
que son muy útiles en muchos problemas.
Sin ir más allá miremos la figura:

\begin{figure}[H]
	\begin{center}
		\begin{asy}
	pair A = dir(80);
	pair B = dir(160);
	pair C = dir(210);
	pair D = dir(330);

	pair X = 1.3 * (B - C) + C;
	pair Y = 1.2 * (A - D) + D;
	draw(X--B);
	draw(A--Y);

	draw(A--B--C--D--cycle);

	draw(A--C);
	draw(B--D);

	draw(unitcircle);

	dot("$A$", A, NE);
	dot("$B$", B, NW);
	dot("$C$", C, SW);
	dot("$D$", D, SE);

	markangle(A, B, X, radius=10, blue);
	markangle(A, D, C, radius=10, blue);
	markangle(D, C, B, radius=10, red);
	markangle(Y, A, B, radius=10, red);

	markangle(C, B, D, radius=13, green);
	markangle(C, A, D, radius=13, green);
\end{asy}

		\caption{Cuadrilátero cíclico.}
	\end{center}
\end{figure}

Las igualdades de los ángulos verdes
son por arco capaz
y las otras dos son debidas a que
$\angle CBA + \angle ADC = 180^\circ$.
Más aún, si $P$ es la intersección de las
diagonales, entonces los triángulos
$PBC$ y $PAD$ son opuestamente semejantes
al igual que $PBA$ y $PCD$.

Ahora vemos el resultado más potente de todos:

\begin{lemma}
	Sea $ABCD$ un cuadrilátero convexo.
	Entonces $ABCD$ es cíclico si y solo si
	\[\angle DCA = \angle DBA\]
	o
	\[\angle CBA + \angle ADC = 180^\circ.\]
\end{lemma}

Esto nos dice que, con que tengamos una
sola de esas relaciones, entonces el cuadrilátero
ya es cíclico, y si en su lugar tuviéramos
$\angle ACB = \angle ADB$ o $\angle BAD + \angle DCB = 180^\circ$
entonces el cuadrilátero igualmente sería cíclico
ya que es cuestión de aplicar el teorema
con un orden distinto de los vértices.

Pero nos queda una pregunta: sabemos que
si el cuadrilátero es cíclico, entonces se dan
las igualdades de ángulos, pero si se dan las
igualdades de ángulos, ¿por qué el cuadrilátero
debe ser cíclico?

\centering
\begin{asy}
	pair A = dir(80);
	pair B = dir(160);
	pair C = dir(210);
	pair D = dir(330);

	pair X = 1.3 * (B - C) + C;
	draw(X--B);

	draw(A--B--C--D--cycle);

	draw(A--C);
	draw(B--D);

	draw(unitcircle);

	dot("$A$", A, NE);
	dot("$B^\prime$", B, NW);
	dot("$C$", C, SW);
	dot("$D$", D, SE);
	dot("$B$", X, W);
	draw(X--D);

	markangle(C, B, D, radius=10, blue);
	markangle(C, A, D, radius=10, blue);
	markangle(C, X, D, radius=10, blue);
\end{asy}
\par
\raggedright

Este problema puede ser resuelto planteándolo
por contradicción.
Sea $\Omega$ el circuncírculo de $ADC$.
Supongamos que $\angle CBD = \angle CAD$
pero que $B$ \textbf{no} está en
el circuncírculo de $ADC$.
Sea $B^\prime$ el segundo punto
de intersección de $BC$ con $\Omega$.
Claramente
\[\angle CBD = \angle CAD = \angle CB^\prime D,\]
por arco capaz,
pero esto implica que $DB$ y $DB^\prime$
son rectas paralelas, lo que no es cierto ya
que intersecan en $D$.

El caso $\angle ADC + \angle CBA = 180^\circ$
se deja como un ejercicio al lector.

Ilustramos el uso de este hecho con
el siguiente teorema:

\begin{theorem}
	Las alturas de un triángulo intersecan
	en un punto, a este punto se le denomina
	el ortocentro del triángulo.
\end{theorem}

\begin{figure}[H]
	\begin{center}
		\begin{asy}
	pair A = dir(110);
	pair B = dir(210);
	pair C = dir(330);

	pair D = foot(A, B, C);
	pair E = foot(B, A, C);
	pair F = foot(C, A, B);

	pair H = extension(A, D, B, E);
	draw(A--B--C--cycle);
	draw(B--E--F--C);
	draw(A--H);
	draw(H--D, dashed);

	dot("$A$", A, N);
	dot("$B$", B, SW);
	dot("$C$", C, SE);
	dot("$H$", H, (-2, -0.4));
	dot("$D$", D, S);
	dot("$E$", E, (1.3, 0.3));
	dot("$F$", F, W);
\end{asy}

	\end{center}
	\caption{El ortocentro del triángulo $ABC$.}
\end{figure}

Sean $E$ y $F$ los pies de las
alturas del triángulo desde $B$ y $C$
respectivamente y sea $H$ la intersección de
$BE$ y $CF$. Sea $E$ la intersección
de $AH$ y $BC$.

Notemos que si demostramos que $AD$
es perpendicular a $BC$, entonces el problema
ya estaría resuelto ya que habríamos probado
que las tres alturas coinciden en un punto.

Recordamos nuestro lema de cuadriláteros cíclicos
y vemos que $\angle BFC = \angle BEC = 90^\circ$,
lo que implica que $BFEC$ es cíclico, más aún,
como $\angle AEH + \angle HFA = 90^\circ + 90^\circ = 180^\circ$
entonces $AFHE$ también es cíclico.

Ahora es cuestión de ver igualdades de ángulos:
\[90 - \angle B = \angle FCB = \angle FEB = \angle FEH = \angle FAH,\]
todo porque $BFEC$ y $AFHE$ son cíclicos.
Para finalizar vemos que
\[\angle ADB = 180^\circ - (\angle BAD + \angle B) = 180^\circ - (90 - \angle B + \angle B) = 90.\]

\section{Problemas}

\begin{problem}[Teorema de Pick o Parcelarias]
Se tiene un polígono cuyos lados no se intersecan
en el plano cartesiano
teniendo todos  sus vértices coordenadas enteras.
Si $I$ es la cantidad de puntos
con coordenadas enteras en el interior del polígono
y $B$ la cantidad
en el borde, demuestre que el área del polígono
es igual a
\[I + \frac{B}2 - 1.\]
\end{problem}

\begin{problem}
Sean $O$ y $H$ el circuncentro y ortocentro respectivamente de
un triángulo acutángulo $ABC$.
Demuestre que $\angle BAO = \angle HAC$.
\end{problem}

\begin{problem}[Teorema de Brahmagupta nunca supe cómo se escribe]
Sea $ABCD$ un cuadrilátero cíclico. Se sabe que
sus diagonales $AC$ y $BD$ son perpendidulares. Demuestre que la recta
perpendicular a $BC$ que pasa por la intersección de las diagonales biseca
al lado $AD$.
\end{problem}

\begin{problem}[Recta de Simson]
Un punto $X$ se encuentra en el circuncírculo de un triángulo $ABC$. Los pies de las perpendiculares
desde $X$ a los lados $BC$, $AC$ y $AB$ son $P$, $Q$ y $R$ respectivamente. Demuestre que $P$, $Q$ y $R$
son colineales.
\end{problem}

\begin{problem}
Sean $A$, $B$ y $C$ tres puntos en una recta en
se orden. Sea dibujan las circunferencias
$\Omega$ y $\Gamma$ con diámetros $AB$ y $AC$
respectivamente. Sea $D$ un punto en $\Omega$
tal que $CD$ es tangente a $\Omega$.
La recta $CD$
en $E$. Si $AB = BC$, determine el valor de
\[\frac{CD}{DE}.\]
\end{problem}

\begin{problem}
Sea $ABCD$ un paralelogramo tal que $\angle DBA = 2\angle DBC$. Sea $T$ la intersección
de la recta perpendicular a $BC$ que pasa por $A$ con el segmento $BD$. Si
$AB = 6$, calcule la longitud del segmento $TD$.
\addhints{ricardojaja1}
\end{problem}

\begin{problem}
Sean $a$, $b$, $c$ y $d$ números reales positivos con
\begin{align*}
	a^2 + b^2      & =       c^2 + d^2  \\
	a^2 + d^2 - ad & =  b^2 + c^2 + bc.
\end{align*}
Encuentre el valor de $\frac{ad + bc}{ab + cd}$.
\addhints{ciclicoalgebra1, ciclicoalgebra2, ciclicoalgebra3, ciclicoalgebra4}
\end{problem}

\begin{problem}
Sea $ABCD$ un cuadrilátero convexo con
$AC = AD$ y $\angle BAC = \angle CAD$.
Sea $E$ un punto en $AC$ talque $AE = AB$.
Muestre que $BC = DE$.
\addhints{rotomoteciacuadrilatero}
\end{problem}

\begin{problem}
Sea $ABC$ un triángulo inscrito en una circunferencia $\Omega$ con
$AC$ diámetro de $\Omega$. Sea $P$ la intersección de la recta $AC$ con
la recta tangente a $\Omega$ que pasa por $B$. Si $\angle BAC = 33^\circ$,
halle el valor del ángulo $\angle CPB$.
\end{problem}

\begin{problem}
Sea $ABC$ un triángulo isósceles con $AB = AC$. Sean $P$ y $Q$ puntos
en $AB$ y $AC$ respectivamente tales que $AP = QC$. Sea $O$
la intersección de la altura $A$ del triángulo $ABC$ con el
circuncírculo del triángulo $APQ$. Demuetre que $O$ es el
circuncentro de $ABC$.
\end{problem}

\begin{problem}
Sea $ABC$ un triangulo acutángulo, sea $BE$ la prolongación del lado $AB$ y $CF$
la prolongación del lado $AC$, si $AB=CF$ y $AC=BE$, demuestre que el punto
de intersección entre el circuncírculo del triángulo $ABC$ y la altura
del triangulo $AEF$ es circuncentro de $AEF$.
\end{problem}

\begin{problem}
Sea $ABCD$ un paralelogramo. Sean $E$ y $F$ puntos externos al
paralelogramo tales que los triángulos $AED$ y $CDF$ son equiláteros.
Demuestre que $EBF$ es equilátero.
\end{problem}

\begin{problem}
Sea $ABC$ un triángulo. Sea $P \neq C$ un punto en la recta $BC$ tal que $BP = BC$.
Sea $Q$ el punto medio de $AB$ y sea $R$ la intersección de las rectas $PQ$ y $AC$.
Si $G$ es el baricentro de $ABC$, demuestre que $RG$ es paralela a $AB$.
\end{problem}

\begin{problem}[CAM 2025]
Sea $ABC$ un triángulo acutángulo y $X$ un punto en
el lado $AC$ tal que $CX = 2AX$. Sea $P$ la intersección
de la recta $BX$ con la mediana $A$ de $ABC$. Si el
área de $ABC$ es $60$, calcule el área del triángulo $APX$.
\end{problem}

\begin{problem}
Sea $ABC$ un triángulo acutángulo con ortocentro $H$, sea $P$ la reflexión
de $H$ sobre la recta $BC$ y $F$ el pie de la altura desde $C$. La
circunferencia circunscripta al triángulo $AFP$ intersecta
a la recta $BC$ en dos puntos distintos $X$ e $Y$.
Demuestre que $C$ es el punto medio del segmento $XY$.
\end{problem}

\begin{problem}
Sea $ABC$ un triángulo equilátero de lado $\ell$ y sea $P$ un punto interior
a $ABC$. Sean $P_1$, $P_2$ y $P_3$ las proyecciones de $P$ sobre
los lados $AB$, $AC$ y $BC$ respectivamente. En términos de
$\ell$, calcule la suma
\[PP_1 + PP_2 + PP_3.\]
\end{problem}

\begin{problem}
Si el área de un triángulo es $50$, y dos de sus lados miden
$5$ y $6$, calcule todos los valores posibles del lado restante.
\end{problem}

\begin{problem}
Sea $ABCD$ un cuadrilátero convexo.
Demuestre que los cuatro puntos de intersección
entre las bisectrices de los ángulos
$A$ y $C$ con las de los ángulos $B$ y $D$
se encuentran sobre una circunferencia.
\end{problem}

\begin{problem}[Teorema de la mariposa]
A través del punto medio $C$ de una cuerda arbitraria $AB$ de una circunferencia,
se han trazado dos cuerdas $KL$ y $MN$
($K$ y $M$ se encuentran de un mismo lado de $AB$), $Q$ es el punto de intersección
de $AB$ y $KN$, $P$ es el punto de intersección de $AB$ y $ML$. Demuestre que $QC = CP$.
\end{problem}

\begin{problem}[Libro ruso de geometría ``de primaria'']
En una circunferencia se tienen dos puntos fijos
$A$ y $B$ y un punto móvil $M$. En la prolongación del segmento
$AM$ fuera de la circunferencia se traza el segmento $MN = MB$.
Halle el lugar geométrico de los puntos $N$.
\end{problem}

\begin{problem}
Sea $I$ el incentro de un triángulo $ABC$.
$AI$ intersecta a $BC$ en $D$. Demuestre que
\[\frac{AI}{ID} = \frac{AB + AC}{BC}.\]
\end{problem}

\begin{problem}
Se dibujan dos rectas paralelas $\ell_1$
y $\ell_2$ y un punto
$O$ entre ellas. Por $O$ se traza una recta
secante que corta a las rectas paralelas
en $A_1$ y $A_2$ respectivamente. Sea $X$
un punto tal que $OA_1 = XA_2$ y $XA_2 \perp A_1A_2$.
Encuentre el espacio geométrico de todos los puntos $X$
a medida que la recta secante varía.
\end{problem}

\begin{problem}
Demuestre que en un trángulo
un bisector más corto correspode
a un lado más largo.
\end{problem}

\begin{problem}
Sean $a$ y $b$ dos reales positivos.
Se ubican dos puntos $A$
y $B$ en una recta $\ell$ y los puntos $A_1$
y $B_1$ tal que $AA_1 = a$ y $BB_1 = b$.
con $AA_1 \perp \ell$ y $BB_1 \perp \ell$.
Sea $H$ el punto de intersección de las rectas
$AB_1$ y $A_1B$. Demuestre que la distancia
de $H$ a la recta $\ell$ es constante
sin importar la elección de $A$ y $B$.
\end{problem}

\begin{problem}[Iberoamericana 2024]
Algunos puntos del plano Euclídeo son pintados
del color rojo con la siguiente regla:
si dos vértices de un triángulo con ángulos
$30^\circ$, $60^\circ$ y $90^\circ$ son rojos,
entonces también lo es el tercero.

Demuestre que si los vértices de un
triángulo cualquiera son rojos,
entonces el baricentro de ese triángulo
también es rojo.
\end{problem}
